\documentclass[a4paper,11pt]{article}

% ==========================================
% 1. 宏包配置 (全部移到这里)
% ==========================================
\usepackage[margin=1in]{geometry}
\usepackage{times}
\usepackage{graphicx}
\usepackage{amsmath}
\usepackage{xcolor}
\usepackage[ruled,vlined,linesnumbered]{algorithm2e}
\usepackage{caption}
% 修复 Unicode 符号问题的包
\usepackage{pifont} 

% --- TikZ 绘图设置 ---
\usepackage{tikz}
\usetikzlibrary{shapes, arrows.meta, positioning, calc, shadows, trees, decorations.pathreplacing, patterns}

% ==========================================
% 2. 自定义设置
% ==========================================
% 颜色定义
\definecolor{sedarRed}{RGB}{220, 53, 69}
\definecolor{morphGreen}{RGB}{25, 135, 84}
\definecolor{codeBg}{RGB}{248, 249, 250}
\definecolor{noteBg}{RGB}{255, 248, 225}
\definecolor{procOrange}{RGB}{253, 126, 20}

% 方便引用的命令
\newcommand{\Morph}{\textsc{Morph}}
\renewcommand{\algorithmicrequire}{\textbf{Input:}}
\renewcommand{\algorithmicensure}{\textbf{Output:}}

% 定义带圈数字 (替代 Unicode 字符)
\newcommand{\cnum}[1]{\ding{\numexpr171+#1}}

\title{Debug: Syntax-Aware Query Abstraction}
\author{Morph Team}
\date{\today}

\begin{document}

\maketitle

\section{Approach}

\subsection{Syntax-Aware Query Abstraction}
\label{sec:abstraction}

The primary objective of this stage is to decouple the \textit{logical intent} of a test case from its \textit{dialect-specific representation}. To achieve this, we introduce a core data structure termed the \textbf{Abstract Query Template (AQT)}.

\paragraph{Definition 3.1 (Abstract Query Template).} An AQT is a generalized Abstract Syntax Tree (AST) $\mathcal{T}_{abs}$ derived from a source query $Q_{src}$. In $\mathcal{T}_{abs}$, dialect-specific nodes (e.g., proprietary functions $f_{spec}$) are replaced by high-level abstract tokens $\tau \in \Sigma_{abs}$ (e.g., \texttt{<TimeDiff>}), while the invariant relational algebra structure (e.g., \texttt{JOIN}, \texttt{WHERE}) and data flow dependencies are strictly preserved.

\subsubsection{Abstraction Process}
As depicted in Figure~\ref{fig:aqt_trans}, the abstraction process transforms a concrete PostgreSQL AST into a dialect-agnostic AQT. Unlike simple regex replacement, this is a structural analysis governed by our static analysis engine.

% ============================================================
% Figure: Abstraction Engine Visualization
% ============================================================
\begin{figure}[t]
\centering
\begin{tikzpicture}[
    level distance=1.2cm,
    sibling distance=2.2cm,
    font=\footnotesize\ttfamily,
    edge from parent/.style={draw, -latex, gray, thick},
    % --- 样式定义 ---
    concrete/.style={rectangle, draw=blue!50, fill=blue!5, rounded corners, minimum height=0.6cm, align=center},
    abstract/.style={rectangle, draw=morphGreen!50, fill=morphGreen!5, rounded corners, dashed, minimum height=0.6cm, text=morphGreen, font=\bfseries},
    boundary/.style={rectangle, draw=sedarRed!80, fill=sedarRed!10, rounded corners, thick, minimum height=0.6cm, align=center, drop shadow={opacity=0.2, color=sedarRed}},
    engineBox/.style={rectangle, draw=procOrange!60, top color=procOrange!10, bottom color=procOrange!20, rounded corners=6pt, font=\sffamily\bfseries\small, align=center, minimum height=1.8cm, minimum width=2.8cm, drop shadow={opacity=0.2, color=procOrange}, inner sep=8pt},
    discarded/.style={rectangle, draw=gray!30, fill=gray!10, rounded corners=2pt, font=\tiny\ttfamily, text=gray!60, inner sep=2pt},
    dialogBox/.style={rectangle, draw=gray!40, fill=noteBg, rounded corners=3pt, font=\scriptsize\sffamily, align=left, inner sep=5pt, drop shadow={opacity=0.1}, text width=3.8cm},
    % 修复箭头参数错误: radius -> width/length
    pointer/.style={-{Circle[open, width=3pt, length=3pt]}, thick, gray!50, shorten >=2pt, dashed}
]

% 1. Left Tree: Source AST
\begin{scope}[local bounding box=srcTree]
    \node[concrete] (root) {SELECT}
        child { node[concrete] {WHERE}
            child[sibling distance=3cm] { node[concrete] (op) {>}
                child { node[concrete] {u.created\_at} }
                child { 
                    node[boundary] (subRoot) {(SUBQUERY\\ROOT)} 
                    child[level distance=1cm] { node[concrete, font=\scriptsize, fill=white, opacity=0.7] (subContent) {SELECT MIN(...) ...} }
                }
            }
        };
    \node[above=0.3cm of root, font=\sffamily\bfseries] {Source AST (Concrete)};
\end{scope}

% 2. Right Tree: AQT Skeleton (加大间距到13cm)
\begin{scope}[xshift=13cm, local bounding box=tgtTree]
    \node[concrete] (root2) {SELECT}
        child { node[concrete] {WHERE}
            child[sibling distance=3cm] { node[concrete] {>}
                child { node[concrete] {u.created\_at} }
                child { node[abstract] (abstractToken) {<SubQuery>} }
            }
        };
    \node[above=0.3cm of root2, font=\sffamily\bfseries] {AQT Skeleton (Normalized)};
\end{scope}

% 3. Middle: Abstraction Engine (绝对坐标定位)
\path let \p1 = (root), \p2 = (root2) in 
    node[engineBox] (engine) at ({(\x1+\x2)/2}, -3.5) {
    Abstraction\\Engine
    \vspace{0.2em}
    \begin{tikzpicture}[scale=0.25, baseline=-3pt, rotate=15]
        \foreach \a in {0,45,...,315} \fill[procOrange!70] (\a:1) -- (\a+12:1.2) -- (\a+24:1) -- cycle;
        \fill[procOrange!70] (0,0) circle (1); \fill[white] (0,0) circle (0.4);
    \end{tikzpicture}
};

% Arrows
\draw[-{Latex[width=3mm]}, line width=1.5pt, gray!50] (srcTree.east |- engine) -- (engine.west);
\draw[-{Latex[width=3mm]}, line width=1.5pt, morphGreen!60] (engine.east) -- (tgtTree.west |- engine) 
    node[midway, above=-2pt, font=\scriptsize\sffamily\bfseries, text=morphGreen] {Extract Structure};
\draw[-{Latex[width=2mm]}, line width=1pt, gray!40, dashed] (engine.south) -- ++(0, -0.8) 
    node[midway, right, font=\scriptsize\sffamily, text=gray] {Discard Details};

% Discarded bits
\node[discarded, below=0.2cm of engine.south, xshift=-0.5cm, rotate=-10] {SELECT MIN...};
\node[discarded, below=0.5cm of engine.south, xshift=0.4cm, rotate=5] {::INTERVAL};
\node[discarded, below=0.35cm of engine.south, xshift=0cm, rotate=15] {AGE()};

% 4. Explanation Dialogs
% 修复 Unicode 字符，使用 \cnum{}
\node[dialogBox, below=2.8cm of srcTree.south west, anchor=north west, xshift=0.5cm] (step1) {
    \textbf{\textcolor{sedarRed}{\cnum{1} Boundary Detection}}\\
    Traverser hits a structural boundary (e.g., subquery root). The algorithm recognizes this requires special handling.
};
\draw[pointer, sedarRed] (step1.north) to[out=90, in=220] (subRoot.south west);

\node[dialogBox, right=0.5cm of step1] (step2) {
    \textbf{\cnum{2} Process \& Filter}\\
    The engine separates invariant structure from dialect-specific details. Concrete subtrees are pruned and discarded.
};
\draw[pointer, procOrange] (step2.north) to[out=90, in=250] (engine.south west);
\node[font=\Huge, text=sedarRed, opacity=0.8] at (subContent) {$\times$};

\node[dialogBox, right=0.5cm of step2] (step3) {
    \textbf{\textcolor{morphGreen}{\cnum{3} Normalized Result}}\\
    A generic placeholder (\texttt{<SubQuery>}) is injected. The skeleton preserves the logical "shape" for reuse.
};
\draw[pointer, morphGreen] (step3.north) to[out=90, in=320] (abstractToken.south east);

\end{tikzpicture}
\caption{\textbf{Visualization of the Abstraction Engine.} The engine identifies complex structural boundaries (\cnum{1}), filters out dialect-specific subtrees (\cnum{2}), and injects generic tokens to produce a normalized skeleton (\cnum{3}).}
\label{fig:aqt_trans}
\end{figure}

The extraction logic is formalized in \textbf{Algorithm~\ref{alg:aqt_extraction}}. The algorithm operates in two phases to transform the raw AST into a structured AQT tuple.

\textbf{Phase 1: Structure Decomposition (Lines 7-15).}
The function \textsc{Traverse} recursively explores the AST. Unlike traditional flat translation, it employs a structure-aware strategy: whenever it encounters a structural boundary (e.g., a subquery or CTE, checked by \textsc{IsBoundary}), it halts the linear traversal. Instead of processing the internal details of the subquery, it creates a new segment edge in the Structure Graph $\mathcal{G}$ and replaces the entire subtree with a structural tag \texttt{<SubQuery>} (Line 10).

\textbf{Phase 2: Dialect Normalization (Lines 11-13).}
For non-boundary nodes, the algorithm consults the rule set $\mathcal{R}$. If a node matches a dialect-specific feature (e.g., `AGE()`), it is abstracted into a dialect-agnostic token (e.g., `<TimeDiff>`). Standard SQL keywords remain unchanged. This ensures the resulting $S_{norm}$ retains the precise logical "shape" without carrying over the syntactic noise of the source dialect.

% ============================================================
% Algorithm 1: Structure-Aware AQT Generation (Compact)
% ============================================================
% \vspace{-0.5cm}
\begin{algorithm}[ht]
\small
\caption{Structure-Aware AQT Generation}
\label{alg:aqt_extraction}

\SetKwInOut{Input}{Input}
\SetKwInOut{Output}{Output}
\SetKwProg{Fn}{Function}{:}{}

\Input{Source query $Q_{src}$; Schema $\mathcal{S}$; Abstraction rules $\mathcal{R}$}
\Output{AQT $\mathcal{T} = \langle S_{norm}, \mathcal{D}, \mathcal{G}, \mathcal{H}, \mathcal{C} \rangle$}

$\mathcal{G} \leftarrow \emptyset$; $Segs \leftarrow \emptyset$ \tcp*{Structure graph; segment storage}
$AST \leftarrow \textsc{Parse}(Q_{src}, \mathcal{S})$ \tcp*{Build AST with schema binding} \label{line:alg1:parse}
$root \leftarrow \textsc{Traverse}(AST.root, \textsc{Null})$ \tcp*{Recursive decomposition}
$S_{norm} \leftarrow \textsc{Normalize}(Segs[root])$ \tcp*{Normalized skeleton} \label{line:alg1:norm}
$\mathcal{H} \leftarrow \textsc{SHA256}(S_{norm})$; $\mathcal{C} \leftarrow \textsc{Complexity}(AST)$ \tcp*{Hash; score}
$\mathcal{D} \leftarrow \textsc{ExtractSchema}(\mathcal{S}, Q_{src})$ \tcp*{Schema context}
\Return $\langle S_{norm}, \mathcal{D}, \mathcal{G}, \mathcal{H}, \mathcal{C} \rangle$\; \label{line:alg1:return}

\hrulefill

\Fn{\textsc{Traverse}($node$, $pid$)}{ \label{line:alg1:traverse_start}
    $id \leftarrow \textsc{NewID}()$; $tokens \leftarrow [\,]$\;
    \lIf{$pid \neq \textsc{Null}$}{$\mathcal{G}.\textsc{AddEdge}(pid \to id)$} \label{line:alg1:addedge}
    \ForEach{$c$ \textbf{in} $node.children$}{
        \uIf{\textsc{IsBoundary}($c$)}{ \label{line:alg1:boundary_start}
            $cid \leftarrow \textsc{Traverse}(c, id)$ \tcp*{Recurse into subquery/CTE}
            $tokens.\textsc{Append}(\texttt{<SubQuery:}cid\texttt{>})$\; \label{line:alg1:boundary_end}
        }
        \uElseIf{$r \leftarrow \textsc{MatchRule}(c, \mathcal{R})$}{ \label{line:alg1:rule_start}
            $tokens.\textsc{Append}(r.abstract)$ \tcp*{e.g., AGE()$\to$<TimeDiff>} \label{line:alg1:rule_end}
        }
        \lElse{$tokens.\textsc{Append}(c.text)$} \label{line:alg1:else}
    }
    $Segs[id] \leftarrow tokens$; \Return $id$\; \label{line:alg1:segs}
} \label{line:alg1:traverse_end}
\end{algorithm}

\end{document}